\section*{Exercise 4}
\textit{Task: Plotting a coin flip and evaluating the results} \newline

\noindent In this task, we evaluate and visualize the outcomes of repeated coin flips to explore the usefulness of plotting and to identify potential problems in interpreting visual data.

\noindent Our approach was to generate coin flip data using a \textbf{Java} program, save the results in a CSV file, and then use \textbf{Python} with the \textbf{pandas} and \textbf{matplotlib} libraries for visualization. In our encoding scheme, we represent heads as \texttt{1} and tails as \texttt{-1}. We created arrays of length $n = 1000$ to store the individual outcomes of the flips.

\noindent To develop an initial understanding of the data, we first plotted the absolute counts of heads and tails, as shown in Figure~\ref{fig1}. In addition, we created a cumulative sum plot to observe the long-term trend in the results.

\begin{figure}[h]
    \centering
    \includegraphics[width = 7cm]{assignment2/plots/absflip.png}
    \hfill
    \includegraphics[width = 7.4cm]{assignment2/plots/cumflip.png}
    \caption{Coin flip results for $n=1000$ flips: (left) absolute count of heads vs. tails; (right) cumulative sum over time.}
    \label{fig1}
\end{figure}

\noindent The plots illustrate clearly that the expected equal distribution between heads and tails does not always hold, especially in smaller sample sizes. For instance, in Figure~\ref{fig1}, the number of tails exceeds the number of heads — despite the theoretical probability being 50/50. This discrepancy becomes apparent in both plots.
This is also visible in the mean and the variance of the plots: 
$$
\text{mean:}\  \overline{x} = -0.074 
$$$$ \text{variance:}\  s^2_n = 0.995$$

\noindent The second plot (relative frequencies) offers an additional insight: although local variability are present, the overall trend appears to stabilize toward an equal ratio as the number of trials increases. This highlights an important concept in probability, showing that outcomes tend to converge to the expected distribution as the number of trials grows\newline

\noindent To answer the core question of Exercise 4: What happens if we start with \( x = 0 \), then toss a coin \( N \) times, adding 1 for heads and subtracting 1 for tails?

\noindent Figure~\ref{fig:my_label} visualizes exactly this process. At each step, the cumulative sum \( x \) is updated depending on whether the coin flip resulted in heads (+1) or tails (-1). After 1000 tosses, the graph ends at around \( x = -32 \), which at first glance might suggest that tails occurred significantly more often than heads.

\noindent However, this impression is misleading. The final relative deviation is just $\frac{-32}{1000} = -0.032$ meaning that tails occurred only 3.2\% more often than heads. Over the course of 1000 trials, this is well within the expected randomness.

\begin{figure}[h]
    \centering
    \includegraphics[width = \textwidth]{assignment2/plots/coinflipwalk.png}
    \caption{"Coin Flip Walk": Evolution of cumulative sum \( x \) when adding 1 for heads and subtracting 1 for tails}
    \label{fig:my_label}
\end{figure}

\noindent The plot reveals an important characteristic of a simple random walk: although the individual steps are unbiased (equal probability for heads and tails), the cumulative sum can drift away from zero for long stretches. In this particular simulation, the walk tends to stay below zero, giving the visual impression that tails dominated, even though the actual imbalance is small.

\noindent Unlike relative frequency plots, this kind of walk does not stabilize around a particular value, because each outcome affects all future values. If we increased the number of tosses further, we would expect the walk to continue fluctuating around zero, though not necessarily converging back to it.

\noindent Repeating the experiment multiple times shows that in some cases the walk drifts positively, in others negatively — confirming that the apparent trend in a single run should not be overinterpreted.\newline



\noindent Another important visualization is shown in Figure~\ref{fig2}, where we compare the expected value with the cumulative sum over time for five independent simulations.

\begin{figure}[h]
    \centering
    \includegraphics[width = 1.0\textwidth]{assignment2/plots/5sum.png}
    \caption{Comparison of expectation and simulation over time for five independent runs}
    \label{fig2}
\end{figure}

\noindent In Figure~\ref{fig2}, we can see that as the number of coin flips increases, the results become more stable and get closer to what we expect on average. This shows the Law of Large Numbers, which means that the more times we repeat the experiment, the closer the average result gets to the true probability..

\noindent Additionally, this plot highlights the practical importance of using a grid in visualizations. The absence of grid lines in this figure makes it difficult to estimate exact values or interpret the convergence behavior at specific points in time. Proper grid usage significantly improves the readability and interpretability of such plots.\newline

\newpage\noindent Figure~\ref{fig4} shows the results for 100{,}000 coin tosses. These plots illustrate how, despite the randomness of each individual toss, the overall trends begin to stabilize when the number of trials increases significantly. 

\noindent The first two plots show the absolute counts of heads and tails as well as their relative frequencies, both clearly converging toward equality as expected from probability theory. The third plot visualizes the cumulative sum (adding +1 for heads and -1 for tails), which still shows some fluctuation, but remains within a narrower range relative to the number of tosses.

\noindent This again confirms the nature of random walks: even with 100,000 tosses, randomness causes small local deviations, but the global behavior tends toward balance.

\begin{figure}[h]
    \centering
    \includegraphics[width=7cm]{assignment2/plots/100k.1.png}
    \includegraphics[width=7cm]{assignment2/plots/100k.2.png}
    \includegraphics[width=15cm]{assignment2/plots/100k.3.png}
    \caption{Results of 100,000 coin tosses: (1) relative frequencies, (2) absolute values, and (3) cumulative sum over tosses.}
    \label{fig4}
\end{figure}
